Met de opkomst van cloud computing en het Internet of Things is ook de
architectuur van moderne applicaties ingrijpend veranderd. Waar deze vroeger
bestonden uit monolithische programma's die op lokale systemen in eigen beheer
draaiden, is er nu een duidelijke verschuiving naar een architectuur gebaseerd
op kleinere, onderling verbonden \emph{services}. Deze worden uitgerold op
infrastructuur die gedeeld wordt met meerdere gebruikers. Het delen van
resources met andere applicaties biedt namelijk diverse voordelen, zoals
verbeterde schaalbaarheid, hogere efficiëntie en lagere kosten. Daartegenover
staan echter ook nieuwe uitdagingen, met name op het gebied van cybersecurity.

Voor een correcte werking van een applicatie is het essentieel om de
confidentialiteit, integriteit en authenticiteit van gegevens en programmacode
te waarborgen. Hoewel er diverse technieken bestaan voor veilige opslag en
verzending van gegevens, blijft de bescherming van gegevens terwijl ermee
gerekend wordt een actief onderzoeksdomein. \acp{TEE} trachten dit probleem op
te lossen door gebruik te maken van hardware-gebaseerde mechanismen. Deze
mechanismen beschermen zowel de programmacode als de gegevens van een programma
in uitvoering, of deze nu aanwezig zijn in de CPU of in het geheugen, en
schermen deze dus af van andere programma's die binnen hetzelfde systeem worden
uitgevoerd. \acp{TEE} bieden zelfs bescherming tegen geprivilegieerde software,
zoals de hypervisor of het besturingssysteem, waardoor het risico op
kwetsbaarheden aanzienlijk wordt verkleind.

In dit proefschrift analyseren we \acp{TEE} vooral vanuit een praktisch
perspectief en onderzoeken we hoe deze veilig en effectief kunnen worden
toegepast in zowel bestaande als nieuwe toepassingen, rekening houdend met de
beperkingen van deze technologie. Allereerst stellen we een kader voor dat de
ontwikkeling en het beheer van gedistribueerde toepassingen in gedeelde,
heterogene omgevingen vereenvoudigt. Hiervoor maken we gebruik van
procesgebaseerde \acp{TEE}, zoals Intel \ac{SGX} en Sancus, die de attack
surface aanzienlijk verkleinen en veilige I/O mogelijk maken. Dit biedt sterke
end-to-end integriteitsgaranties. Om de effectiviteit van ons werk te
demonstreren, ontwikkelen en evalueren we een prototype dat gebaseerd is op een
smart home systeem.

In een tweede bijdrage bestuderen we meer recente architecturen, zoals AMD
\sevsnp{} en Intel \ac{TDX}, die het mogelijk maken volledige virtuele machines
binnen een \ac{TEE} uit te voeren. Hoewel deze technologieën het ontwikkelen van
toepassingen voor een \ac{TEE} vereenvoudigen, vergroten ze zowel de attack
surface als de operationele complexiteit. Een virtuele machine bestaat immers
uit meerdere softwarecomponenten die tijdens het opstartproces worden uitgevoerd
en het is essentieel is om elke fase van dit proces zorgvuldig te verifiëren.
Onze bevindingen tonen echter aan dat gebruikers in huidige cloudomgevingen
beperkte kennis hebben van deze software. Bovendien hebben cloudproviders vaak
aanzienlijke controle over hoe virtuele machines worden opgestart en kunnen ze
mogelijk zelfs toegang krijgen tot gevoelige gebruikersgegevens die door de
\ac{TEE} worden beschermd. Het is daarom van cruciaal belang om de
vertrouwensrelatie tussen cloudproviders en klanten te duidelijk te stellen bij
het gebruik van \acp{TEE} in de cloud.

Als derde bijdrage onderzoeken we de interacties tussen een applicatie die
binnen een \ac{TEE} wordt uitgevoerd en het onderliggende, mogelijk
gecompromitteerde systeem. We bestuderen in het bijzonder hoe een stabiel en
betrouwbaar tijdsbegrip binnen een \ac{TEE} kan worden gerealiseerd. Daarbij
beoordelen we verschillende oplossingen en bespreken we veelvoorkomende
toepassingen die afhankelijk zijn van tijd voor een correcte werking. Op basis
van deze analyse ontwikkelen we een innovatieve methode voor de revocatie van
digitale certificaten in \ac{V2X}-toepassingen, waarbij we \acp{TEE} gebruiken
in de voertuigen. We tonen formeel aan dat ons voorstel een bovengrens voor de
revocatietijd garandeert. Onze evaluatie bevestigt bovendien de efficiëntie en
schaalbaarheid van deze aanpak.

Samengevat, dit proefschrift toont aan hoe \acp{TEE} kunnen worden gebruikt om
enerzijds de beveiliging van bestaande applicaties te verbeteren en anderzijds
om nieuwe, innovatieve oplossingen te ontwikkelen die bijkomende voordelen
bieden in vergelijking met de meer traditionele oplossingen. Daarnaast hebben we
ook aandacht voor de praktische uitdagingen die ontstaan bij het implementeren
van \ac{TEE}-toepassingen in minder veilige omgevingen zoals publieke clouds.
